\chapter{Valoraciones y propuestas de mejora}
\label{ch:opinion_conclusión}

\section{Papel del tutor de empresa}
\label{sec:rolTutor}
Describir el papel que ha tomado el tutor de empresa, Santiago Rodríguez Palma, durante las prácticas resulta complicado debido a que las prácticas han sido una formación en la que el papel de tutor iba rotando semanalmente o cada pocos días.

Por lo tanto, se va a realizar una opinión sobre el rol que han tomado de forma general todos los tutores a lo largo de los 2 meses de formación.

Todos los tutores que han impartido la formación han aportado una gran cantidad de materia teórica y, especialmente, experiencia en sus respectivos departamentos, dado que todos ellos eran expertos con varios años de experiencia trabajando en la empresa.

En todos los casos, se ha permitido e incentivado interactuar durante las lecciones, preguntando dudas o pidiendo desarrollar un tema en más profundidad. Además, como se ha explicado con anterioridad, cada lección cuenta con su fase práctica y de evaluación y, en ambos casos, se ha resuelto cada duda o problema que surgiese.


\section{Valoración}
\label{sec:valoracion}
En lo personal, todos los tutores me han parecido excelentes en su abarcamiento de las competencias impartidas. Todo el contenido dado se me ha explicado de forma correcta y se me han dado las herramientas, materiales y entornos necesarios para poner a prueba los conceptos aprendidos.

\section{Propuestas de mejora}
\label{sec:improvements}

Lo único negativo que podría mencionar de estas prácticas sería la mala organización del tiempo.

Entre las 4 horas diarias entre lunes y jueves, las tareas diarias, los cursos Udemy y la presentación final, las horas dedicadas a esta formación superan significativamente lo planteado inicialmente en el anexo. No considero que en sí sea un problema. Cada hora que se ha invertido ha resultado productiva y valen el esfuerzo que requieren. Sin embargo, creo que sería relevante indicar previamente la expectativa de horas reales y, especialmente, considerar ajustar la fecha en las que se realizan.

La combinación de no indicar de forma correcta la duración de las practicas y su solapamiento con la fase de exámenes ordinaria y extraordinaria de la UCLM causan una sobrecarga de trabajo que pueden llevar a percibir negativamente las prácticas y repercutir en los resultados académicos.

